% from ega2-1-fr.tex
\oldpage[II]{41}
\begin{lemma}[2.7.9.1]
\label{II.2.7.9.1-fr}
Sous les hypothèses de (\hyperref[II.2.7.5-fr]{2.7.5}), pour tout $n\geq0$,
il existe un entier $k$ (dépendant de $n$) et un homomorphisme surjectif
$\mathcal{O}_X^k\to\mathcal{O}_X(n)$.
\end{lemma}

Il suffit en effet (\hyperref[II.2.7.2-fr]{2.7.2}) de montrer que, pour un $k$
convenable, il y a un homomorphisme (TN)-surjectif $u$ \emph{de degré $0$},
du $S$-module gradué produit $S^k$ dans $S(n)$. Or, on a $(S(n))_0=S_n$, et
par hypothèse $S_h=S_1^h$ pour tout $h>0$, donc
$SS_n=S_n+S_{n+1}+\cdots$. Comme $S_n$ est un $S_0$-module de type fini
(\hyperref[II.2.1.5-fr]{2.1.5} et
\hyperref[II.2.1.6-fr]{2.1.6, (i)}), considérons un système de générateurs
$(a_i)_{1\leq i\leq k}$ de ce module~; l'homomorphisme $u$ fera correspondre
$a_i$ au $i$-ème élément de la base canonique de $S^k$ ($1\leq i\leq k$)~;
comme $\operatorname{Coker}u$ s'identifie alors à
$(S(n))_{-n}+\cdots+(S(n))_{-1}$, $u$ répond bien à la question.
\end{proof}

\begin{corollary}[2.7.10]
\label{II.2.7.10-fr}
Sous les hypothèses de (\hyperref[II.2.7.5-fr]{2.7.5}), soit $\mathcal{F}$ un
$\mathcal{O}_X$-Module quasi-cohérent de type fini. Il existe un entier $n_0$
tel que, pour tout $n\geq n_0$, $\mathcal{F}$ soit isomorphe à un quotient
d'un $\mathcal{O}_X$-Module de la forme $(\mathcal{O}_X(-n))^k$ ($k$
dépendant de $n$).
\end{corollary}

\begin{proposition}[2.7.11]
\label{II.2.7.11-fr}
Supposons vérifiées les hypothèses de
(\hyperref[II.2.7.5-fr]{2.7.5}) et soit $M$ un $S$-module gradué. Alors~:
\begin{enumerate}
  \item[\rm(i)] L'homomorphisme canonique
    $\widetilde{\alpha}:\widetilde{M}\to
    \widetilde{\Gamma_*(\widetilde{M})}$ est un isomorphisme.
  \item[\rm(ii)] Soit $\mathcal{G}$ un sous-$\mathcal{O}_X$-Module
    quasi-cohérent de $\widetilde{M}$, et soit $N$ le sous-$S$-module gradué
    de $M$, image réciproque de $\Gamma_*(\mathcal{G})$ par $\alpha$. Alors
    on a $\widetilde{N}=\mathcal{G}$ ($\widetilde{N}$ étant identifié, en
    vertu de (\hyperref[II.2.5.4-fr]{2.5.4}), à un
    sous-$\mathcal{O}_X$-Module de $\widetilde{M}$).
\end{enumerate}
\end{proposition}

\begin{proof}
Comme $\beta:\widetilde{\Gamma_*(\widetilde{M})}
\to\widetilde{M}$ est un isomorphisme en vertu de
(\hyperref[II.2.7.5-fr]{2.7.5}), $\widetilde{\alpha}$ est l'isomorphisme
réciproque en vertu de (\hyperref[II.2.6.5.1-fr]{2.6.5.1}), d'où (i). Soit
$P$ le sous-module gradué $\alpha(M)$ de $\Gamma_*(\widetilde{M})$~; comme
$M\mapsto\widetilde{M}$ est un foncteur exact
(\hyperref[II.2.5.4-fr]{2.5.4}), l'image de $\widetilde{M}$ par
$\widetilde{\alpha}$ est égale à $\widetilde{P}$, donc, en vertu de (i),
$\widetilde{P}=\widetilde{\Gamma_*(\widetilde{M})}$. Posons
$Q=\Gamma_*(\mathcal{G})\cap P$~; en vertu de ce qui précède et de
(\hyperref[II.2.5.4-fr]{2.5.4}), on a
$\widetilde{Q}=\widetilde{\Gamma_*(\mathcal{G})}$, donc la
restriction de $\beta$ à $\widetilde{Q}$ est un \emph{isomorphisme} de ce
$\mathcal{O}_X$-Module sur $\mathcal{G}$ par
(\hyperref[II.2.7.5-fr]{2.7.5}). Mais par définition de $N$ et
(\hyperref[II.2.5.4-fr]{2.5.4}), la restriction de l'isomorphisme
$\widetilde{\alpha}$ à $\widetilde{N}$ est un isomorphisme de $\widetilde{N}$
sur $\widetilde{Q}$, d'où la conclusion par
(\hyperref[II.2.6.5.1-fr]{2.6.5.1}).
\end{proof}

\subsection{Comportements fonctoriels.}
\label{subsection:II.2.8-fr}

\begin{env}[2.8.1]
\label{II.2.8.1-fr}
Soient $S$, $S'$ deux anneaux gradués à degrés positifs,
$\varphi:S'\to S$ un homomorphisme d'anneaux gradués. Nous désignerons par
$G(\varphi)$ la partie ouverte de $X=\operatorname{Proj}(S)$ complémentaire
de $V_+(\varphi(S'_+))$, ou, ce qui revient au même, la réunion des
$D_+(\varphi(f'))$, où $f'$ parcourt l'ensemble des éléments homogènes de
$S'_+$. La restriction à $G(\varphi)$ de l'application continue
${}^a\varphi$ de $\operatorname{Spec}(S)$ dans $\operatorname{Spec}(S')$
(\hyperref[I.1.2.1-fr]{I, 1.2.1}) est donc une application continue de
$G(\varphi)$ dans $\operatorname{Proj}(S')$, que nous noterons encore
${}^a\varphi$ par abus de langage. Si $f'\in S'_+$ est homogène, on a
\[
\label{II.2.8.1.1-fr}
  {}^a\varphi^{-1}(D_+(f'))=D_+(\varphi(f'))
\tag{2.8.1.1}
\]
compte tenu du fait que ${}^a\varphi$ applique $G(\varphi)$ dans
$\operatorname{Proj}(S')$ et de
(\hyperref[I.1.2.2.2-fr]{I, 1.2.2.2}). D'autre part, l'homomorphisme
$\varphi$ définit canoniquement (avec les mêmes notations) un homomorphisme
d'anneaux gradués $S'_{f'}\to S_f$, d'où, par restriction aux éléments de
degré $0$,

